\documentclass[
	12pt,						 % Schriftgröße
	DIV=10,						 % Satzspiegelgröße Seitenränder
	german,						 % für Umlaute, Silbentrennung etc.
	paper=a4,					 % Papierformat
	oneside, 
	titlepage,					 % es wird eine Titelseite verwendet
	parskip=half,				 % Abstand zwischen Absätzen (halbe Zeile)
	headings=normal,			 % Größe der Überschriften verkleinern
	listof=totoc,				 % Verzeichnisse im Inhaltsverzeichnis aufführen
	bibliography=totoc,			 % Literaturverzeichnis im Inhaltsverzeichnis aufführen
	numbers=noenddot,			 % Kapitelnummern ohne Punkt	
	enabledeprecatedfontcommands,% Alte Befehle zur unterstüzung von \listofsymbols	
    notitlepage
]{scrreprt}

\makeatletter	
	\setlength\@fptop{0\p@}  % Entfernt Lücken am Seitenbeginn
\makeatother

\newcommand{\titel}{Entwicklung einer Mikrocontroller-Basierten Lösung zum Auslesen analoger Zähler mittels KI-Bilderkennung}
\newcommand{\art}{Praxistransferbericht}
\newcommand{\autor}{Claas Meints}
\newcommand{\matrikel}{172489}
\newcommand{\semester}{4}
\newcommand{\studienbereich}{Elektrotechnik/Mechatronik}
\newcommand{\studiengang}{Elektrotechnik}
\newcommand{\erstgutachter}{Prof. Dr. Herwig Henseler}
\newcommand{\ewe}{EWE AG}
\newcommand{\geraet}{SiGI-Device}
\newcommand{\comment}[1]{}
\newcommand{\TODO}{}
\newcommand\fig[2]{
	\begin{figure}[ht]
		\centering
		\includegraphics[width=0.8\textwidth]{figures/#1}
		\caption{#2} 
		\label{fig:#1}
	\end{figure}
}
\newcommand{\ESPTDM}{\cite{ESP32}\cite{ESP32-S2}\cite{ESP32-C3}\cite{ESP32-S3}}
% Zur Richtigen Verwendung von Einheiten
\usepackage[
    locale=DE, 
    %per=symbol-or-fraction
]{siunitx}

\sisetup{
    mode = text,
    detect-family,
    detect-weight,  
    exponent-product = \cdot,
    output-decimal-marker={\text{.}},
    %math-rm=\mathsf,
    %text-rm=\sffamily,
    range-phrase = { -- },
    per-mode=symbol,
    per-symbol=/
}
\DeclareSIUnit{\var}{var}

% Anpassung des Seitenlayouts -------------------------------------------------------------% 	siehe Seitenstil.tex
% -------------------------------------------------------------
\usepackage[headsepline,automark]{scrlayer-scrpage}

\usepackage{listings}
\usepackage{color}
% C++ color scheme
\definecolor{backgroundcolor}{RGB}{237, 241, 237}
\definecolor{commentcolor}{RGB}{0,128,0}
\definecolor{keywordcolor}{RGB}{136, 134, 58}
\definecolor{identifiercolor}{RGB}{24, 79, 110}
\definecolor{numbercolor}{RGB}{128,128,128}
\definecolor{preprocessorcolor}{RGB}{128,128,128}
\definecolor{stringcolor}{RGB}{163,21,21}

\lstdefinestyle{C++}{
    backgroundcolor=\color{backgroundcolor},   
    commentstyle=\color{commentcolor},
    directivestyle=\color{preprocessorcolor},
    keywordstyle=\color{keywordcolor},
    identifierstyle=\color{identifiercolor},
    numberstyle=\tiny\color{numbercolor},
    stringstyle=\color{stringcolor},
    basicstyle=\ttfamily\footnotesize,
    breakatwhitespace=false,         
    breaklines=true,                 
    captionpos=b,                    
    keepspaces=true,                 
    numbers=left,                    
    numbersep=5pt,                  
    showspaces=false,                
    showstringspaces=false,
    showtabs=false,                  
    tabsize=2
}
\lstset{language=C++,escapechar=@,style=C++}
\usepackage{booktabs}
\usepackage{eurosym}
\usepackage{pifont}% http://ctan.org/pkg/pifont

\renewcommand{\lstlistingname}{Code} % Damit Codebeispiele mit Code betitelt werden.


% Anpassung an Landessprache -------------------------------------------------------------% 	Verwendet globale Option german siehe \documentclass
% -------------------------------------------------------------
\usepackage[ngerman]{babel}
\usepackage{scrhack}
\usepackage[utf8]{inputenc}
\usepackage{amsmath,amssymb,units}
\usepackage{wrapfig,caption, placeins}
\usepackage{multicol}
\captionsetup{
    format = plain,
    justification = centering,
    labelsep = newline,
    singlelinecheck = false,
    labelfont = bf,
    font = small
}
\usepackage{mathptmx,charter,helvet,courier}
\usepackage[printonlyused, smaller]{acronym}
\usepackage{wrapfig}
\usepackage{pgfplots}
% Umlaute -------------------------------------------------------------% 		Umlaute/Sonderzeichen wie äöüß direkt im Quelltext verwenden (CodePage).
%		Erlaubt automatische Trennung von Worten mit Umlauten.
% -------------------------------------------------------------
\usepackage[utf8]{inputenc}
\usepackage[T1]{fontenc}
\usepackage{ae}			 	% "schöneres" ä
\usepackage{textcomp} 		% Euro-Zeichen etc.
\usepackage[german=quotes]{csquotes}


% Verwendung von Blindtext zur Layout gestaltung -------------------------------------------------------------
\usepackage{blindtext}

% Grafiken -------------------------------------------------------------
% 		Einbinden von Grafiken [draft oder final]
% 		Option [draft] bindet Bilder nicht ein - auch globale Option
% -------------------------------------------------------------
\usepackage{graphicx}
%\graphicspath{{Grafiken/}} 	% Dort liegen die Bilder des Dokuments


% Befehle aus AMSTeX für mathematische Symbole z.B. \boldsymbol \mathbb ----
\usepackage{amsmath,amsfonts}
\usepackage{tensor}

% Für Index-Ausgabe; \printindex -------------------------------------------------------------
\usepackage{makeidx}

% Einfache Definition der Zeilenabstände und Seitenränder -------------------------------------------------------------
\usepackage{setspace}
\usepackage{geometry}


%% Zum Umfließen von Bildern -------------------------------------------------------------
\usepackage{floatflt}


% Lange URLs umbrechen etc. -------------------------------------------------------------
\usepackage{url}

% für lange Tabellen
\usepackage{longtable}
\usepackage{multirow}
\usepackage{rotating}
\usepackage{diagbox}
\usepackage{array}
\newcolumntype{M}[1]{>{\centering\arraybackslash}m{#1}}
\usepackage{ragged2e}
\usepackage{lscape}
\usepackage{colortbl} %farbige hinterlegung


% Spaltendefinition rechtsbündig mit definierter Breite auf dem Deckblatt 
%------------------------------------------------------------
\newcolumntype{d}[1]{>{\raggedleft\hspace{0pt}}p{#1}}

% Formatierung von Listen ändern
\usepackage{paralist}
% Standardeinstellungen:
\setdefaultleftmargin{2.5em}{2.2em}{1.87em}{1.7em}{1em}{1em}

%Zur Verwendung von BibLaTex
\usepackage[
    style=alphabetic
]{biblatex}
\addbibresource{bibliografie.bib}
\DefineBibliographyStrings{ngerman}{
    andothers = {{et\,al\adddot}},             
}

\usepackage{xcolor}
\usepackage{caption}
\usepackage{float}
\usepackage{subfig}
\usepackage{xspace}
\usepackage{scrhack}

\usepackage[acronyms,toc]{glossaries} %Glossar Achtung Perl.exe muss installiert werden
\usepackage{glossaries-extra}
\usepackage{glossary-mcols}

% pdf-Optionen  -------------------------------------------------------------
\usepackage{pgfplots} % Für Kompatibilät unter verschiedenen OS
\pgfplotsset{compat=1.7}

\pdfminorversion=7 % Verwende pdf Version 1.7

\usepackage[
    bookmarks,
    bookmarksopen=true,
    bookmarksnumbered=true,
    colorlinks=true,
    %linkcolor=red, % einfache interne Verknüpfungen
    %anchorcolor=black,% Ankertext
    %citecolor=blue, % Verweise auf Literaturverzeichniseinträge im Text
    %filecolor=magenta, % Verknüpfungen, die lokale Dateien öffnen
    %menucolor=red, % Acrobat-Menüpunkte
    %urlcolor=cyan, 
    % für die Druckversion können die Farben ausgeschaltet werden:
    linkcolor=black, % einfache interne Verknüpfungen
    anchorcolor=black,% Ankertext
    citecolor=black, % Verweise auf Literaturverzeichniseinträge im Text
    filecolor=black, % Verknüpfungen, die lokale Dateien öffnen
    menucolor=black, % Acrobat-Menüpunkte
    urlcolor=black, 
    %backref,
    %pagebackref,
    plainpages=false,% zur korrekten Erstellung der Bookmarks
    pdfpagelabels,% zur korrekten Erstellung der Bookmarks
    hypertexnames=false,% zur korrekten Erstellung der Bookmarks
    linktoc=all%Sowohl Seitenzahl als auch Text als Link
]{hyperref}
\usepackage{bookmark}
\usepackage[all]{hypcap} % jump to top of figure
\usepackage{pdfpages}
% Zeilenabstand ------------------------------------------------------------
\onehalfspacing


% Seitenränder -------------------------------------------------------------
\geometry{paper=a4paper,left=25mm,right=25mm,top=30mm,bottom=40mm}

% Kopf- und Fußzeilen ------------------------------------------------------
\pagestyle{scrheadings}

% Kopf- und Fußzeile auch auf Kapitelanfangsseiten -------------------------
\renewcommand*{\chapterpagestyle}{scrheadings}

% Schriftform der Kopfzeile ------------------------------------------------
\renewcommand{\headfont}{\normalfont}

% Kopfzeile ----------------------------------------------------------------
\ihead{
	\parbox{11.8cm}{\normalsize{ \ }}\\
	%\small{\untertitel}\\[2ex]
	\textit{ \headmark}
}
\chead{
		%\includegraphics[height=0.8cm]{Grafiken/EWENETZ_Logo.png}
}
\ohead{	
		\includegraphics[height=1.4cm]{figures/Logo_PHWT.png}	
}
\setlength{\headheight}{21mm} % Höhe der Kopfzeile


% Fußzeile -----------------------------------------------------------------
\ifoot{
	%\includegraphics[height=1.5cm]{Grafiken/Logo_PHWT.png}
}
\cfoot{}
\ofoot{\pagemark}
\setlength{\footheight}{15mm}


% erzeugt ein wenig mehr Platz hinter einem Punkt --------------------------
%\frenchspacing 

% Schusterjungen und Hurenkinder vermeiden
\clubpenalty = 10000
\widowpenalty = 10000 
\displaywidowpenalty = 10000


%% Fußnoten fortlaufend durchnummerieren ------------------------------------
%\counterwithout{footnote}{chapter}
\setkomafont{section}{\large} 



\newglossary[llg]{glossary}{llo}{lls}{Glossar}
\newglossaryentry{Cloud Computing}
{
    name=Cloud Computing,
    description={},
    type=glossary
}
\newglossary[slg]{symbolslist}{syi}{syg}{Symbolverzeichnis}
\newglossaryentry{x}{
    name={\ensuremath{x}},
    description={Variable},
    type=symbolslist
}
\newglossary[wlg]{wordslist}{woi}{wog}{Fremdwortverzeichnis}

\makeglossaries
\newacronym{aws}{AWS}{Amazon Web Services}
\newacronym{iot}{IoT}{Internet of Things}


\begin{document}
\pagenumbering{Roman}
\thispagestyle{plain}
\begin{titlepage}
    \begin{center}
        \huge{\textbf{\art}}\\[1.5ex]
        \LARGE{\titel}\\[2ex]
        \large{Private Hochschule für Wirtschaft und Technik Vechta/Diepholz}
        \begin{tabular}{r l}
            vorgelegt von: & \quad \autor \\
        \end{tabular}
    \end{center}
    \setcounter{tocdepth}{0}\let\clearpage\relax\tableofcontents
    {\centering\textcopyright\ 2022\\[1.5ex]}
    \singlespacing\small\noindent Dieses Werk einschließlich seiner Teile ist \textbf{urheberrechtlich geschützt}. Jede Verwertung außerhalb der engen Grenzen des Urheberrechtgesetzes ist ohne Zustimmung des Autors unzulässig und strafbar. Das gilt insbesondere für Vervielfältigungen, Übersetzungen, Mikroverfilmungen sowie die Einspeicherung und Verarbeitung in elektronischen Systemen. 
\end{titlepage}

\chapter{Einleitung}\label{ch:1}
\pagenumbering{arabic}

Der Begriff des Cloud Computings genießt aktuell einen hohen Stellenwert in der Informationstechnologie und verändert die IT Landschaft maßgeblich. Seit mehr als 10 Jahren schreitet das Anbieten von IT-Dienstleistungen in der Wolke kontinuierlich voran und die absolute Mehrheit an IT-Unternehmen weltweit nutzt Cloud Computing Angebote oder denkt über diese nach\cite*[]{Cloud}. Insbesondere für kleine und mittelständische Unternehmen bietet das Auslagern von Rechenleistung und weiteren Diensten interessante Möglichkeiten. Denn die Anschaffung und Wartung eines eigenen Rechenzentrums stellt kleinere Unternehmen vor eine Herausforderung\cite*[]{DAAS}. Doch nicht nur für Unternehmen bietet die Cloud neue Möglichkeiten, sondern auch Privatleute profitieren von den Vorteilen der Cloud. Die Meisten Privatpersonen verbinden Cloud bisher mit der Sicherung von ihren Daten auf einem Server.

Neben der Cloud existiert aktuell ein weiterer Trend der ebenfalls eine hohe Aufmerksamkeit genießt, nämlich das Internet of Things (zu deutsch: Internet der Dinge). Dieser Begriff beschreibt die fortschreitene Vernetzung von Produkten, Services, Maschinen und Sensoren mittels IP-basierter Protokolle. Ähnlich wie die Cloud spielt das IoT eine wichtige Rolle für die Industrie (smart factories, predictive maintenance) aber auch private Haushalte (smart Homes)\cite*[]{Next}\cite*[]{Öster}. 

Beide Themen sind für sich schon Trends die aktuell viel Aufmerksamkeit genießen. Umso interessanter ist die Kombination der beiden Themen. In einem IoT Netzwerk können riesige Mengen an Daten anfallen, welche nicht immer Zentral auf einem IoT Gerät verarbeitet werden können. Daher bietet sich das Verarbeiten der Daten in einer Cloud an. In der Cloud werden nun auf Grundlage der gesammelten Daten Berechnungen durchgeführt und beispielsweise mithilfe Künstlicher Intelligenz Aktionen bestimmt, die anschließend ausgeführt werden. 

Thema dieser Ausarbeitung ist eben genau dieses Zusammenspiel der beiden Themenfelder. Hierzu wird zuerst auf die Grundlagen des Cloud Computings eingegangen und anschließend das Themenfeld IoT beleuchtet. Nach den Grundlagen wird kurz auf Amazon Web Services eingegangen und im Anschluss ein Beispielprojekt vorgestellt, in dem IoT Geräte mit AWS verknüpft werden.  

\chapter{Was ist \Gls{Cloud Computing}}\label{ch:2}
%\fig{2.0}{IoT}
\section{Geschichte und Urpsrung}
Bei der Cloud handelt es sich nicht um eine neue Technologie sondern eher um eine Wiederverwendung mehrerer bekannter Technologien. Um also das Konzept hinter Cloud Computing zu verstehen wird kurz áuf die einzelnen Vorgänger eingegangen.

Dir Grundidee reicht bereits in die 1960er Jahre zurück. Bereits damals bestand der Grundgedanke, größere Probleme auf mehrere Rechner zu verteilen und somit effizienter zu berechnen. Das hierbei enstehende System nennt sich verteiltes System und die einzelnen Rechner werden Knoten genannt. Dies wurde damls in Form von "Cluster Computing" durchgeführt. Hierbei werden identische Rechner durch ein hochgeschwindigkeits Netzwerk miteinander Verbunden, was zu einer völlig neuen Rechenleistung und Verfügbarkeit führte. Diese Lösung stellt sich als deutlich günstiger als ein Großrechner heraus und wird noch heute in Supercomputern angewandt.

Wenn das Netzwerk, anders als beim Grid-Computing, aus Knoten mit verschiedenen Aspekten aufgabeut ist, wird vom Grid-Computing gesprochen. Die verschiednen Knoten können sich hierbei in Hardware, Software oder ihrer Anbindung an das Netzwerk unterscheiden. Zudem könne sich die verschiedenen Knoten vom Netzwerk trennen. Dieses Konzept wird häufig dazu genutzt, um unbenutzte CPUs verschiednener Rechner auszunutzen und somit Große Probleme in kleinen Teilproblemen auf verschiedenen Rechnern zu bearbeiten. 

Ein weiteres Konzept stellt das Utility-Computing dar. Der Grundgedanke hierbei ist es, dass der Nutzer nur seine genutzte Rechenleistung bezahlt. Oft wird hier der Vergleich mit der Steckdose gezogen, da hier der Kunde ebenfalls nur für seinen bezogenen Strom bezahlt[https://d-nb.info/1065664664/34]. Etwaige überlegungen sind jedoch an der nicht verfügbaren Hardware und dem mangehaften Netzausbau gescheitert. Schließlich wurde das Konzept in den 1990er Jharen in das Konzept der Cloud integriert und ist fester Bestandteil des Konzepts.

 
\section{Definition}
Eine einheitliche Definition  für Cloud Computing existiert in der Literatur nicht. Vorallem weil der Begriff relativ unscharf ist und sich die Technologie noch im Wandel befindet. Nach besteht Cloud Computing aus einer Ansammlung von Diensten, Anwendungen oder ITRessourcen, die dem Nutzer flexibel und skalierbar über das Internet angeboten werden. AWS als größter Cloud Anbieter definiert Cloud computing wie folgt: "Cloud Computing" ist die bedarfsabhängige Bereitstellung von IT-Ressourcen über das Internet zu nutzungsabhängigen Preisen. Statt physische Rechenzentren und Server zu erwerben, zu besitzen und zu unterhalten, können Sie über einen Cloud-Anbieter wie Amazon Web Services (AWS) nach Bedarf auf Technologieservices wie beispielsweise Rechenleistung, Speicher und Datenbanken zugreifen." [https://aws.amazon.com/de/what-is-cloud-computing/] 

\section{Merkmale}
Da es jedoch keine einheitliche Definition gibt lässt sich Cloud Computing lediglich auf Grundlage der Merkmale Charakterisieren. Die Merkmale des National Institute of Technologys and Standards (NIST) gelten in der Praxis als auch Forschunge als anerkannt und sind in Abbildung \ref{fig:2.2} zu sehen
\fig{2.2}{https://www.uni-muenster.de/imperia/md/content/angewandteinformatik/aktivitaeten/publikationen/daas-fuer-kmu.pdf}
\subsection{Ressource Pooling}
Der Begriff Ressource Pooling beschreibt das Ansammeln von verschiedensten IT-Ressourcen zu einem Pool auf den die Benutzer zugreifen. Der Benutzer greift aus Seiner Sicht auf Ressourcen wie z.B. eine virtuell Maschine, Speicherkapazität oder eine Dienstsinstanz zu. Mithilfe von Virtualisierung lassen sich physische Ressourcen für den Benutzer auf virteulle Ressourcen abbilden und gleichzeitig besteht eine Trennung. Der Benutzer sieht also nicht was unter seiner virtuellen Maschine physisch liegt. Durch die Virtualisierung lassen sich Rechenzentren erheblich effizienter betreiben. Die Auslastung lässt sich hierbei statt der herkömmlichen 8 - 15 Prozent auf 70 - 80 Prozent steigern. Dies ist Interessant, da Server während sie untätig sind fast dieselbe Leistung benötigen, als wären sie ausgelastet [https://www.vmware.com/content/dam/digitalmarketing/vmware/en/pdf/vmware-reduce-power-consumption-wp.pdf]. Der Benutzer hat somit keinen Einfluss auf die physischen Ressourcen auf welchen seine Dienste ausgeführt werden.Viele Cloud Anbieter bieten jedoch die Option an den Standort der physischen Ressourcen einzugrenzen z.B. aufgrund von Datenschutzrichtlinien.  [Buyya et al. 2008]
\subsection{Rapid Elasticity}
Der Cloud ist es möglich auf Änderungen der Rechenleistung schnell und dynamisch zu reagieren. Die meisten Anbierter können dies auch ohne Vorlaufzeit, sodass der Benutzer das Gefühl suggeriert bekommt, aus unendlichen Ressourcen zu schöpfen. Die Anpassung erfolgt hierbei händisch durch den Benutzer oder automatisch durch den Anbieter. Wenn also z.B. ein Geschäft zu Weihnachten eine höhere Anzahl an Anfragen auf ihrer Website erwartet, passt sich die Cloud automatisch an die steigenden Anfragen an. 
[Vaquero et al. 2009].
\subsection{On-demand Self-service}
Eine weitere Charakteristik der Cloud ist die Möglichkeit sich nach Bedarf selbst bedienen  zu können. Der Cloudanbierter muss somit dem Benutzer auf anfrage die angeforderte Rechenleistung bereitstellen ohne dabei viel Konfiguration vornehmen zu müssen um die Elastizität zu gewährleisten. Der wirtschaftliche Betrieb ist hierbei nur durch automatisierte Verwaltung und Optimierung gewährleistet und bildet ein wichtiges Merkmal eines Cloud Anbieters.
[Wang et al. 2008]
\subsection{Broad Network Access}
In den meisten Fällen erfolgt der Zugriff auf die Ressourcen der Cloud über ein Netzwerk. Um einen Zugriff von diversen Geräten zu ermöglichen, werden standardisierte REST Schnittstellen oder Web-APIs und Protokolle bzw. Datenformate wie z.B. HTTP, XML JSON usw. verwendet. Vorraussetzung hierfür ist jedoch eine ausreichende Bandbreite der Verbindung.
\subsection{Measured Service}

[Vaquero et al. 2009],

\section{Bereitstellungsmodelle}
[Vaquero et al. 2009]

\section{Service-Ebenen}
\fig{2.1}{Wirtz, B.W. (2018), Electronic Business, 6., aktualisierte und erweiterte Auflage 2018, Wiesbaden 2018.}
\subsection{Infrastructure as a Service}
\subsection{Platform as a Service}
\subsection{Software as a Service}

\section{Vor- und Nachteile der Cloud}
\chapter{\acrfull{iot}}\label{ch:3}
\section{Was ist \acrshort{iot}}\label{sec:3.1}
Unter \acrshort{iot} wird die Vernetzung vieler physicher Geräte mittels dem Internet verstanden. Der Begriff ist bereits seit 1999 in Gebrauch\cite*[]{IARJSET} und wird für eine vielzahl unterschiedlicher technischer Lösungen und Protokolle verwendet. Die verschiedenen technischen Lösungen eint dabei die Verwendung des \acrlongpl{ip} (\acrshort{ip}) innerhalb des \Gls{Protokoll-Stack}s\cite*[]{Rayes2019}.\\ Der erste Schritt in Richtung \acrshort{iot} ist dabei für gewöhnlich die Ausstattung von Geräten mit einer vielzahl von Sensoren. So können Geräte den Zustand ihrer Umgebung erfassen und diese über die Netzwerkverbindung an ein \Gls{Backend} übertragen. Dadurch wird es möglich physische Zustände lokal, aber potenziell auch von überall auf der Welt aus zu überwachen. Das Backend kann also zum einen einer Privatperson ermöglichen die eigenen Geräte zu überwachen, zum anderen können die Geräte auch als Datenquellen zum Beispiel für Unternehmen dienen, welche dann die Informationen vieler Geräte in einem \Gls{Cloud}-\Gls{Backend} bündeln\cite*[]{inproceedings}.\\ Für viele technische Anwendungen ist neben Sensordaten auch die Möglichkeit der Steuerung von Aktoren notwendig. Hierbei ist sowohl die Vernetzung mit anderen \acrshort{iot}-Geräten als auch die Vernetzung mit dem \Gls{Backend} wichtig. Über die Kopplung mit einer \acrshort{iot}-Sensoreinheit kann so zum Beispiel die Regelung des Aktors vorgenommen werden. Wird der Aktor an ein Backend angeschlossen, ist darüber eine automatische Fernsteuerung möglich\cite*[]{Rayes2022}.\\ Neben Sensoren und Aktoren sind Einheiten für die Verarbeitung von Daten notwendig. Dabei kann die Datenverarbeitung sowohl im \Gls{Backend} geschehen, als auch durch ein \acrshort{iot}-Gerät. Die Verarbeitungslogik kann sowohl in ein bestehendes \acrshort{iot}-Gerät integriert werden, als auch als separate Einheit ausgeführt sein. Durch die so geschaffene Möglichkeit der lokalen Verarbeitung von Daten wird das Gesamtsystem resilienter gegen den Ausfall von Netzwerkverbindungen. Auch die Latenz wird durch die direktere Verbindung im lokalen Netzwerk verringert\cite*[]{Rayes2022}.
\section{Anwendungsgebiete}\label{sec:3.2}
Durch die bereits beschriebene Vielfältigkeit der technischen Lösungen, die unter dem Begriff \acrshort{iot} zusammengefasst werden ergibt sich auch eine Vielzahl an Anwendungsgebieten und Einsatzmöglichkeiten.\\ Einer der präsentesten Bereiche in dem \acrshort{iot}-Geräte eingesetzt werden ist die Hausautomatisierung. Hier eingesetzte \acrshort{iot}-Geräte können zum einen dem Komfort der Benutzer*innen erhöhen, in dem Arbeiten im Haushalt automatisiert oder erleichtert werden. Darüber hinaus existieren bereits viele Lösungen, welche primär der Sicherheit gegen Einbruch etc. dienen. Eine weitere wichtige Aufgabe von \acrshort{iot}-Geräten im Haus ist die energetische Optimierung. So werden \glqq{}intelligente\grqq\ Heizsysteme mit Temperaturfühlern und Aktoren an den Rolläden gekoppelt, um die gewünschte Temperatur im Innenraum möglichst kostengünstig und Emissionsarm zu erreichen. Auch die Anpassung des elektrischen Energieverbrauchs an die Verfügbarkeit und den Preis von elektrischer Energie kann mittels \acrshort{iot} erfolgen\cite*[]{Zaheeruddin2020}.\\ Auch für die Sammlung medizinisch relevanter Daten werden vermehrt \acrshort{iot}-Geräte eingesetzt. Ein Beispiel kann die Messung von Blutruck oder die Durchfühung eines Elektrokardiogramms durch ein Armband sein. Da diese permanent getragen werden, können Gesundheitliche Probleme auch ohne Verdachtsfall erkannt werden. Eine rechtzeitige Behandlung wird so wahrscheinlicher\cite*[]{Zaheeruddin2020}.\\ In der Landwirtschaft können \acrshort{iot}-Geräte dazu verwendet werden zum Beispiel das Pflanzenwachstum zu überwachen. Durch die Erfassung der Daten kann optimal bewässert und gedüngt werden. Auch vollautomatische Systeme können mittels Steuerungstechnik realisiert werden. \acrshort{iot} bietet dabei die Möglichkeit Sensordaten zu erfassen und im \Gls{Cloud}-\Gls{Backend} mit Wetterdaten zu verknüpfen um langfristig die Bewässerung planen und Optimieren zu können\cite*[]{Zaheeruddin2020}.\\ Der Vermutlich wichtigste Einsatzbereich von \acrshort{iot} ist die Industrie. Der Einsatz von \acrshort{iot}-Geräten in der Industrie wird in Deutschland mit dem Begriff \glqq{}Industrie 4.0\grqq\ beschrieben. Mögliche aufgaben von \acrshort{iot}-Geräten sind die Bereitstellung von Daten zur Bestimmung wartungsintervallen oder aber auch die Automatisierung ganzer Produktionsstraßen. Viele Unternehmen streben dabei eine sogenannte \glqq{}Smart Factory\grqq\ an, also eine Produktionsstraße, welche unterschiedliche Produkte fertigen kann und die Produktionsdaten von einem oder mehreren in dem Produkt integrierten \acrshort{iot}-Geräten erhält. Alle an dem Fertigungsprozess beteiligten Geräte sind dabei untereinander und mit dem Produkt vernetzt. Im idealfall ist für ein Angepasstes Produkt nur eine Änderung der im Produkt gespeicherten Produktionsdaten notwendig, ohne dass die am Fertigungsprozess beteiligten Geräte neu Programmiert oder anderweitig angepasst werden müssen\cite*[]{Zaheeruddin2020}\cite*[]{10.1007/978-3-030-39875-0_20}.\\ % TODO: China?
\section{Technologische Umsetzung von \acrshort{iot}}\label{sec:3.3}
Da eine Ethernet-Verbindung für \acrshort{iot}-Geräte oft nicht verfügbar oder für den Anwendungsfall ungeeignet ist, kommunizieren ein Großteil der \acrshort{iot}-Geräte Drahtlos. So können auf den Layern Eins und Zwei zum Beispiel die energiesparenden Funkstandards IEEE 802.15.4 oder LPWAN bzw. LoRaWAN eingesetzt werden. Spielt der Energieverbrauch eines \acrshort{iot}-Geräts keine übergeordnete Rolle, da das Gerät zum Beispiel permanent mit Strom versorgt wird, werden auch WLAN oder der Mobilfunk für die Kommunikation eingesetzt (vgl. \autoref{fig:3.1}).\fig{3.1}{Funkstandards\cite*[]{Rayes20192}} Auf Layer Drei werden für die \Gls{Protokoll-Stack}s hauptsächlich IPv4 und IPv6 bzw. 6LoWPAN verwendet\cite*[]{Rayes20192}.\\ Auf den darüberligenden Layern gibt es eine Vielzahl unterschiedlicher Prokotkolle und Standards mit denen\acrshort{iot}-Geräte kommunizieren. Sie erfüllen jeweils unterschiedliche Anforderungen, wie Latenz oder Sicherheit, die an die Kommunikation gestellt werden. Darüber hinaus verwenden einige Unternehmen eigene proprietäre Standards für die Anbindung ihrer Geräte. Ein Beispiel für ein besonders in der Hausautomatisierung verbreitetes ofennes Protokoll ist MQTT. Neben der Verwendung spezieller \acrshort{iot}-Protokolle gibt es auch die Möglichkeit Websocket- oder REST-APIs für die Kommunikation zu verwenden\cite*[]{Rayes20192}.\FloatBarrier
\chapter{Amazon Web Services}\label{ch:4}
\acrfull{aws}

Das Unternehmen Amazon bietet selbst mehr als Erfolgreich Cloud Computing Dienste unter dem Namen AWS an. Anfänglich wurden nicht genutzte Rechenkapazitäten verkauft, jedoch hat sich das Anbieten von IT-Ressourcen in der Cloud als extrem profitabel gezeigt und mittlerweile erzielt das Cloud Geschäft des Unternehmens mehr Gewinn als der bekannte Online Handel. Bei AWS handelt es sich aktuell um den Cloud Anbieter mit dem weltweit größten Marktanteil und derzeit werden den Benutzern mehr als 200 verschiedene Dienste angeboten.  [https://de.statista.com/infografik/20802/weltweiter-marktanteil-von-cloud-infrastruktur-dienstleistern/
] 
Im Bereich IaaS bietet der Anbieter beispielsweise Services wie EC2,S3 und SimpleDB an. 
\chapter{ESP32 + \acrshort{aws} Core \acrshort{iot}}\label{ch:4}
\chapter{Zusammenfassung}\label{ch:5}

\listoffigures
% \listoftables
\printglossary[type=\acronymtype,title={Abkürzungsverzeichnis}] % list of acronyms
% \printglossary[type=symbolslist, style=long] % list of symbols
% \printglossary[type=wordslist,   style=long] % list of english words
\printglossary[type=glossary,style=mcoltreenoname] % main glossary
\printbibliography

\appendix
\chapter{Anhang}\label{ch:A}
\section{Quizz Buzzer Codebeispiel (gekürzt)}\label{sec:code}
\lstinputlisting[language=C++]{code/subscribe_publish_sample.c}\newpage

\end{document}